\chapter{生体暗号の解錠：メタファーによる総当たり探索（ブルートフォース）}
\section{天才たちの探索戦略：局所精密走査と超並列ハッキング}
\subsection{ミルトン・エリクソン：診察室における逐次型マイクロ・ブルートフォース}
\subsection{リチャード・バンドラー：多人数空間を一網打尽にする並列ブルートフォース}
\subsubsection{フェイルセーフとエンタメの為のユーモアとそのモデリング}
催眠技法
レートリミットを超える
鍵開けの為のブルートフォースと辞書攻撃
フェイルセーフとしてのジョーク
ミルトン・エリクソンがゆっくり話、バンドラーがマシンガントークで沢山のジョークを喋るわけ。
エリクソンは、診察室で事前アセスメント等でS/N比の高いエビデンスを集めることができた。
エリクソンは鍵の候補が少ない。
バンドラーは、ステージでエリクソンより得られるエビデンスの質が少ないのと、大量の人間を扱うから。
バンドラーは鍵の候補があるから時間内に詰め込まないといけない。


